\chapter{Results}

This chapter presents the empirical findings of the analysis. Each figure is interpreted in detail.

\section{Descriptive Statistics}

Table~\ref{tab:descriptive_ch4} presents descriptive statistics for the analytical sample.

\begin{table}[htbp]
	\centering
	\caption{Descriptive Statistics by Survey Year}
	\label{tab:descriptive_ch4}
	\begin{tabular}{@{}lccc@{}}
		\toprule
		\textbf{Characteristic} & \textbf{KMIS 2015} & \textbf{KMIS 2020} & \textbf{Combined} \\
		\midrule
		Number of children & 3,059 & 3,485 & 6,544 \\
		Number of clusters & 224 & 296 & 520 \\
		Malaria prevalence (\%) & 5.1 & 4.5 & 4.8 \\
		ITN use rate (\%) & 58.0 & 50.5 & 54.1 \\
		Urban population (\%) & 30.5 & 27.5 & 29.0 \\
		\bottomrule
	\end{tabular}
\end{table}

Several patterns emerge from Table~\ref{tab:descriptive_ch4}. First, malaria prevalence declined slightly from 5.1\% in 2015 to 4.5\% in 2020. Second, ITN use decreased from 58.0\% to 50.5\% over the same period. This decline is concerning because it occurred despite continued distribution campaigns, suggesting that distribution alone is insufficient to maintain high usage. Third, the urban population remained stable at approximately 29\% of the sample.

\subsection{Malaria Prevalence by Wealth Quintile}

Figure~\ref{fig:malaria_by_wealth} shows malaria prevalence stratified by wealth quintile for both survey years.

\begin{figure}[htbp]
	\centering
	\includegraphics[width=0.8\textwidth]{Figures/malaria_by_wealth.png}
	\caption{Malaria Prevalence by Wealth Quintile (KMIS 2015 vs 2020)}
	\label{fig:malaria_by_wealth}
\end{figure}

\textbf{Interpretation of Figure~\ref{fig:malaria_by_wealth}:}

This figure displays malaria prevalence on the y-axis (percentage of children testing positive) and wealth quintile on the x-axis, from Poorest (1) to Richest (5). Two lines are shown: blue for 2015 and red for 2020.

Key observations:

\begin{enumerate}
	\item \textbf{Strong wealth gradient}: The poorest children have approximately seven times higher malaria risk than the richest children. In 2015, prevalence was 5.4\% among the poorest compared to 0.8\% among the richest. In 2020, the gradient persists: 4.5\% among the poorest versus 0.7\% among the richest.
	
	\item \textbf{Declining prevalence across all quintiles}: Both surveys show lower prevalence in 2020 compared to 2015 for every wealth quintile. The largest absolute decline occurred in the poorest quintile (0.9 percentage points), while the smallest decline occurred in the richest quintile (0.1 percentage points).
	
	\item \textbf{Non-linear relationship}: The relationship between wealth and malaria is not linear. The largest drop in prevalence occurs between Poorest and Poorer quintiles. The gradient flattens among the top three quintiles.
\end{enumerate}

\textbf{Implications for causal analysis:} Wealth must be carefully controlled in all models. The strong gradient indicates that wealth is a powerful confounder: wealthier households have both higher ITN use and lower baseline malaria risk. Failing to control for wealth would bias estimates toward finding a harmful effect of ITNs (since ITN users are wealthier and wealthier children have lower malaria risk).

\subsection{ITN Use by Wealth Quintile}

Figure~\ref{fig:itn_by_wealth} shows ITN use rates stratified by wealth quintile.

\begin{figure}[htbp]
	\centering
	\includegraphics[width=0.8\textwidth]{Figures/itn_by_wealth.png}
	\caption{ITN Use Rate by Wealth Quintile (KMIS 2015 vs 2020)}
	\label{fig:itn_by_wealth}
\end{figure}

\textbf{Interpretation of Figure~\ref{fig:itn_by_wealth}:}

This figure displays ITN use percentage on the y-axis and wealth quintile on the x-axis.

Key observations:

\begin{enumerate}
	\item \textbf{Positive wealth gradient}: ITN use increases with wealth. In 2015, ITN use ranged from 42.9\% among the poorest to 66.6\% among the richest. In 2020, the range was 32.3\% to 58.4\%.
	
	\item \textbf{Declining use across all quintiles}: Both surveys show lower ITN use in 2020 compared to 2015 for every wealth quintile. The largest absolute decline occurred among the poorest (10.6 percentage points), followed by the poorer (8.4 percentage points). The smallest decline occurred among the richest (8.2 percentage points).
	
	\item \textbf{Persistent gap}: The wealth gap in ITN use remained substantial: approximately 24 percentage points in both years.
\end{enumerate}

\textbf{Combined interpretation of Figures~\ref{fig:malaria_by_wealth} and \ref{fig:itn_by_wealth}:}

Wealthier children have both higher ITN use and lower malaria prevalence. This creates a classic confounding scenario that could produce a biased estimate of ITN effectiveness. If we simply compare ITN users to non-users without adjusting for wealth, we would likely find that ITN users have lower malaria risk, but this would be partly due to wealth rather than ITN use itself. Proper adjustment is essential.

\subsection{Malaria Prevalence by Residence}

Figure~\ref{fig:malaria_by_residence} shows malaria prevalence by urban versus rural residence.

\begin{figure}[htbp]
	\centering
	\includegraphics[width=0.8\textwidth]{Figures/malaria_by_residence.png}
	\caption{Malaria Prevalence by Urban/Rural Residence}
	\label{fig:malaria_by_residence}
\end{figure}

\textbf{Interpretation of Figure~\ref{fig:malaria_by_residence}:}

This bar chart shows malaria prevalence for urban and rural children in 2015 and 2020.

Key observations:

\begin{enumerate}
	\item \textbf{Rural children have 2-3 times higher malaria prevalence}: In 2015, rural prevalence was 6.0\% compared to 2.0\% urban. In 2020, rural prevalence was 3.9\% compared to 1.1\% urban.
	
	\item \textbf{Declining prevalence in both settings}: Both urban and rural areas saw reductions. The absolute reduction was larger in rural areas (2.1 percentage points) than urban areas (0.9 percentage points), but the relative risk remained similar.
\end{enumerate}

\textbf{Implication:} Residence is an important confounder. Rural children have higher baseline malaria risk and may have different ITN access patterns. Residence must be included as a covariate.

\subsection{Malaria Prevalence by Age}

Figure~\ref{fig:malaria_by_age} shows malaria prevalence by child age group.

\begin{figure}[htbp]
	\centering
	\includegraphics[width=0.8\textwidth]{Figures/malaria_by_age.png}
	\caption{Malaria Prevalence by Child Age}
	\label{fig:malaria_by_age}
\end{figure}

\textbf{Interpretation of Figure~\ref{fig:malaria_by_age}:}

This line graph shows malaria prevalence for five age groups: 0-11 months, 12-23 months, 24-35 months, 36-47 months, and 48-59 months.

Key observations:

\begin{enumerate}
	\item \textbf{Lowest among infants (0-11 months)}: Prevalence is approximately 3\% due to maternal antibodies transferred during pregnancy and breastfeeding.
	
	\item \textbf{Peak at 12-23 months}: Prevalence rises to approximately 7\% as maternal antibodies wane and children become more mobile and exposed to mosquitoes.
	
	\item \textbf{Gradual decline after 24 months}: Prevalence gradually declines as children acquire partial immunity through repeated exposure.
	
	\item \textbf{Different peak in 2020}: The 2015 peak occurred at 12-23 months (6.9\%), while the 2020 peak occurred at 36-47 months (4.7\%). This shift may reflect changes in transmission intensity, ITN coverage, or both.
\end{enumerate}

\textbf{Implication:} Age is a strong predictor of malaria risk and must be carefully controlled. The non-linear relationship suggests that linear age terms may be insufficient; flexible modeling (e.g., splines or age groups) is preferable.

\subsection{ITN Use by Age}

Figure~\ref{fig:itn_by_age} shows ITN use by child age group.

\begin{figure}[htbp]
	\centering
	\includegraphics[width=0.8\textwidth]{Figures/itn_by_age.png}
	\caption{ITN Use by Child Age}
	\label{fig:itn_by_age}
\end{figure}

\textbf{Interpretation of Figure~\ref{fig:itn_by_age}:}

This line graph shows ITN use percentage for the same five age groups.

Key observations:

\begin{enumerate}
	\item \textbf{Highest among infants}: ITN use is approximately 66\% (2015) and 52\% (2020) for the youngest children. Caregivers are most diligent about protecting infants.
	
	\item \textbf{Declining with age}: ITN use decreases as children get older. By age 48-59 months, use is approximately 55\% (2015) and 38\% (2020). Older children may resist using nets or may be left uncovered when nets are shared.
	
	\item \textbf{Pattern consistent across both surveys}: The age gradient is stable, though absolute levels are lower in 2020.
\end{enumerate}

\textbf{Combined interpretation of Figures~\ref{fig:malaria_by_age} and \ref{fig:itn_by_age}:}

Younger children have higher ITN use but lower malaria risk (due to maternal antibodies). This creates a situation where ITN users appear to have lower malaria risk partly because they are younger, not solely because of net use. Age adjustment is essential.

\subsection{Positivity Check}

Figure~\ref{fig:overlap} shows the distribution of estimated propensity scores for treated and control groups.

\begin{figure}[htbp]
	\centering
	\includegraphics[width=0.8\textwidth]{Figures/ps_overlap_KMIS_2015.png}
	\caption{Propensity Score Overlap Plot (KMIS 2015)}
	\label{fig:overlap}
\end{figure}

\textbf{Interpretation of Figure~\ref{fig:overlap}:}

This density plot shows the distribution of propensity scores (probability of ITN use given covariates) for treated (red) and control (blue) groups.

Key observations:

\begin{enumerate}
	\item \textbf{Substantial overlap}: The red and blue curves cover largely the same range. For KMIS 2015, treated PS range is 0.058-0.985; control PS range is 0.020-0.940; overlap region is 0.058-0.940.
	
	\item \textbf{No positivity violations}: There are no regions where only treated or only control units exist. For every treated unit, there are comparable control units with similar propensity scores.
	
	\item \textbf{Some treated units have high PS}: Approximately 3.3\% of treated units have PS > 0.95. These are children with very high predicted probability of ITN use. They are successfully matched to control units with PS in the overlap region.
	
	\item \textbf{Some control units have low PS}: Approximately 11.3\% of control units have PS < 0.05. These children have very low predicted probability of ITN use. They are matched to treated units with PS in the overlap region.
\end{enumerate}

\textbf{Conclusion:} The positivity assumption is satisfied, ensuring that causal methods are well-defined. For KMIS 2020, the overlap was similarly good: treated PS range 0.13-0.932, control PS range 0.062-0.856.

\section{Baseline Regression Results}

Table~\ref{tab:baseline_ch4} presents the traditional logistic regression results.

\begin{table}[htbp]
	\centering
	\caption{Baseline Logistic Regression Results (Odds Ratios)}
	\label{tab:baseline_ch4}
	\begin{tabular}{@{}lccc@{}}
		\toprule
		\textbf{Model} & \textbf{KMIS 2015} & \textbf{KMIS 2020} & \textbf{Combined} \\
		\midrule
		Unadjusted & 1.258 (0.896-1.766) & 0.881 (0.640-1.212) & 1.047 (0.831-1.318) \\
		Demo+HH & 1.539 (1.086-2.181) & 1.070 (0.768-1.491) & 1.275 (1.006-1.616) \\
		Fully Adjusted & 0.727 (0.496-1.064) & 0.754 (0.535-1.063) & \textbf{0.777 (0.605-0.999)} \\
		\bottomrule
	\end{tabular}
\end{table}

\textbf{Interpretation of Table~\ref{tab:baseline_ch4}:}

\textit{Unadjusted models:} These models include only ITN use as a predictor. The odds ratios are not statistically significant (all p > 0.05). The point estimates are inconsistent: OR > 1 in 2015 (1.258, suggesting harmful effect) and OR < 1 in 2020 (0.881, suggesting protective effect). This lack of clarity in unadjusted comparisons motivates the need for causal methods.

\textit{Demographic + Household models:} These models add age, sex, wealth, and residence. The odds ratios increase substantially: 1.539 for 2015 (p = 0.015), 1.070 for 2020 (p = 0.688), and 1.275 for combined (p = 0.045). The increase indicates that ITN users are in higher-risk groups conditional on these covariates. After adjusting for demographics, ITN users appear to have higher malaria risk.

\textit{Fully Adjusted models:} These models add all environmental covariates (rainfall, temperature, EVI, aridity, wet days). The coefficients flip direction: all ORs are below 1. For the combined sample, OR = 0.777 (95\% CI: 0.605-0.999, p = 0.049). This demonstrates that environmental factors were masking the true protective effect of ITNs. After controlling for environmental confounders, ITN use shows a statistically significant protective effect.

The pattern across models reveals the direction of confounding:
- Unadjusted: ITN users appear to have higher risk (because they live in high-transmission areas)
- After demography: ITN users appear to have even higher risk (because they are wealthier and wealth is protective)
- After environment: Protective effect emerges (because high-risk areas are controlled for)

\section{GLMM Results}

Table~\ref{tab:glmm_ch4} presents the GLMM results.

\begin{table}[htbp]
	\centering
	\caption{GLMM Results for ITN Use on Malaria}
	\label{tab:glmm_ch4}
	\begin{tabular}{@{}lccc@{}}
		\toprule
		\textbf{Survey} & \textbf{OR} & \textbf{95\% CI} & \textbf{p-value} \\
		\midrule
		KMIS 2015 & 0.557 & (0.352-0.879) & 0.012 \\
		KMIS 2020 & 0.562 & (0.378-0.836) & 0.004 \\
		Combined & 0.588 & (0.436-0.792) & 0.001 \\
		\bottomrule
	\end{tabular}
\end{table}

\textbf{Interpretation of Table~\ref{tab:glmm_ch4}:}

The GLMM estimates are more protective than the fully adjusted logistic regression. ORs range from 0.557 to 0.588, representing 41-44\% reduction in malaria odds within a typical cluster. All are statistically significant (p < 0.05).

\textbf{Why are GLMM ORs smaller (more protective) than fully adjusted logistic ORs?}

GLMM estimates cluster-specific (conditional) effects, holding the random intercept $u_j$ constant. In the presence of unobserved cluster-level heterogeneity, these conditional effects are typically larger in magnitude than marginal (population-average) effects. The logistic regression estimates marginal effects, averaging over cluster-level heterogeneity. With high ICC (50-58\%), the difference is substantial.

Figure~\ref{fig:caterpillar} shows the caterpillar plot of random effects.

\begin{figure}[htbp]
	\centering
	\includegraphics[width=0.8\textwidth]{Figures/glmm_caterpillar_Combined.png}
	\caption{Caterpillar Plot: Cluster Random Effects (Combined Sample)}
	\label{fig:caterpillar}
\end{figure}

\textbf{Interpretation of Figure~\ref{fig:caterpillar}:}

This plot displays the estimated random intercept $u_j$ for each cluster with 95\% confidence intervals.

- \textbf{X-axis}: Clusters sorted by their random effect estimate
- \textbf{Y-axis}: Random intercept $u_j$ (log-odds scale)
- \textbf{Each point}: A cluster's deviation from the overall mean malaria risk
- \textbf{Points above zero}: Clusters with higher-than-average malaria risk
- \textbf{Points below zero}: Clusters with lower-than-average malaria risk
- \textbf{Vertical bars}: 95\% confidence intervals

Key observations:

\begin{enumerate}
	\item \textbf{Substantial variation}: Random intercepts range from approximately -3 to +3 on the log-odds scale. This corresponds to a 20-fold difference in odds between the lowest-risk and highest-risk clusters.
	
	\item \textbf{Many intervals do not cross zero}: Confidence intervals for many clusters do not include zero, indicating statistically significant differences in baseline malaria risk across clusters.
	
	\item \textbf{Consistent with high ICC}: The substantial variation visible in the caterpillar plot is consistent with the ICC of 50-58\%.
\end{enumerate}

\section{Propensity Score Methods with GLMM}

\subsection{PSM Results}

Table~\ref{tab:psm_ch4} presents PSM results with caliper = 0.05.

\begin{table}[htbp]
	\centering
	\caption{PSM Results (Caliper = 0.05)}
	\label{tab:psm_ch4}
	\begin{tabular}{@{}lccc@{}}
		\toprule
		\textbf{Survey} & \textbf{OR} & \textbf{95\% CI} & \textbf{p-value} \\
		\midrule
		KMIS 2015 & 0.392 & (0.194-0.791) & 0.009 \\
		KMIS 2020 & 0.466 & (0.242-0.896) & 0.022 \\
		\bottomrule
	\end{tabular}
\end{table}

\textbf{Interpretation of Table~\ref{tab:psm_ch4}:}

PSM with caliper = 0.05 produces the strongest protective effects among all methods. KMIS 2015 shows a 61\% reduction in odds (OR = 0.392, p = 0.009). KMIS 2020 shows a 53\% reduction (OR = 0.466, p = 0.022). The stronger effect in 2015 may reflect higher ITN use (58.0\% vs 50.5\%) and higher baseline malaria prevalence (5.1\% vs 4.5\%).

Figure~\ref{fig:loveplot} shows the Love plot confirming balance after matching.

\begin{figure}[htbp]
	\centering
	\includegraphics[width=0.8\textwidth]{Figures/psm_love_plot_KMIS_2015_caliper_0.05.png}
	\caption{Love Plot: Balance Assessment (KMIS 2015, Caliper = 0.05)}
	\label{fig:loveplot}
\end{figure}

\textbf{Interpretation of Figure~\ref{fig:loveplot}:}

This Love plot displays standardized mean differences (SMD) for each covariate before matching (blue circles) and after matching (green triangles).

- \textbf{X-axis}: Absolute standardized mean difference (SMD = 0.1 threshold indicated by vertical dashed line)
- \textbf{Y-axis}: Each covariate included in the propensity score model
- \textbf{Good balance}: Points left of the vertical line (SMD < 0.1)

Key observations:

\begin{enumerate}
	\item \textbf{Before matching}: Several covariates show substantial imbalance. Wealth\_numeric has SMD approximately 0.3. Residence has SMD approximately 0.25. Electricity has SMD approximately 0.2. The distance variable (logit PS) has SMD approximately 0.8.
	
	\item \textbf{After matching}: All covariates have SMD < 0.1. The distance variable improved dramatically from SMD $\approx$ 0.8 to SMD $\approx$ 0.03. This indicates that matching successfully balanced all observed confounders.
	
	\item \textbf{The matched treated and control groups are comparable}: After matching, differences in observed covariates are negligible.
\end{enumerate}

\subsection{IPTW Results}

Table~\ref{tab:iptw_ch4} presents IPTW results with truncation at 4.

\begin{table}[htbp]
	\centering
	\caption{IPTW Results (Truncated at 4)}
	\label{tab:iptw_ch4}
	\begin{tabular}{@{}lccc@{}}
		\toprule
		\textbf{Survey} & \textbf{OR} & \textbf{95\% CI} & \textbf{p-value} \\
		\midrule
		KMIS 2015 & 0.482 & (0.319-0.728) & $<$0.001 \\
		KMIS 2020 & 0.511 & (0.349-0.750) & $<$0.001 \\
		\bottomrule
	\end{tabular}
\end{table}

\textbf{Interpretation of Table~\ref{tab:iptw_ch4}:}

IPTW estimates show protective effects of 48-52\% reduction in odds. All are statistically significant at p < 0.001. The estimates are slightly less protective than PSM (which gave 53-61\% reduction). This is expected because PSM discards unmatched units, while IPTW uses the full sample.

Figure~\ref{fig:iptw_weights} shows the IPTW weight distribution.

\begin{figure}[htbp]
	\centering
	\includegraphics[width=0.8\textwidth]{Figures/iptw_weights_KMIS_2015.png}
	\caption{IPTW Weight Distribution (KMIS 2015, Raw vs Truncated)}
	\label{fig:iptw_weights}
\end{figure}

\textbf{Interpretation of Figure~\ref{fig:iptw_weights}:}

The left panel shows raw weights; the right panel shows weights truncated at 4.

- \textbf{X-axis}: Stabilized IPTW weights (should center around 1)
- \textbf{Y-axis}: Count (frequency)
- \textbf{Red bars}: Treated group (ITN users)
- \textbf{Blue bars}: Control group (non-users)

Key observations:

\begin{enumerate}
	\item \textbf{Raw weights}: The distribution centers near 1 for both groups. However, the treated group has some extreme weights up to approximately 9.7 (13 observations with weight > 4). Extreme weights indicate near-positivity violations.
	
	\item \textbf{Truncated weights}: After truncation at 4, the maximum weight is 4. The distribution remains centered near 1. The OR changed minimally from 0.483 to 0.482, confirming that extreme weights did not drive the results.
	
	\item \textbf{KMIS 2020 weights}: The maximum raw weight was 3.89, so no truncation was needed.
\end{enumerate}

\subsection{Doubly Robust Results}

Table~\ref{tab:dr_ch4} presents doubly robust results.

\begin{table}[htbp]
	\centering
	\caption{Doubly Robust Results (Marginal Effects)}
	\label{tab:dr_ch4}
	\begin{tabular}{@{}lccc@{}}
		\toprule
		\textbf{Survey} & \textbf{OR} & \textbf{95\% CI} & \textbf{p-value} \\
		\midrule
		KMIS 2015 & 0.994 & (0.972-1.017) & 0.618 \\
		KMIS 2020 & 0.992 & (0.976-1.009) & 0.370 \\
		\bottomrule
	\end{tabular}
\end{table}

\textbf{Interpretation of Table~\ref{tab:dr_ch4}:}

The doubly robust estimates are near-null (OR $\approx$ 0.99) and not statistically significant. This represents the marginal (population-average) treatment effect.

Why are DR estimates near-null while PSM and IPTW show strong protection? The difference reflects different estimands:

- \textbf{PSM and IPTW (with GLMM)}: Estimate conditional (cluster-specific) effects. They compare treated and untreated children within the same cluster.
- \textbf{DR}: Estimates marginal (population-average) effects using a different estimator that is not conditional on cluster.

When ICC is high (50-58\%), conditional effects are larger than marginal effects. The DR estimate of OR = 0.99 means that, on average across all children in Kenya, ITN use reduces malaria odds by approximately 1\%. This is the population-average effect, which is smaller because it averages over clusters with different baseline risks and different coverage levels.

\section{Double Machine Learning Results}

Table~\ref{tab:dml_ch4} presents DML results using XGBoost.

\begin{table}[htbp]
	\centering
	\caption{DML Results (XGBoost, 5 folds)}
	\label{tab:dml_ch4}
	\begin{tabular}{@{}lccc@{}}
		\toprule
		\textbf{Survey} & \textbf{OR} & \textbf{95\% CI} & \textbf{p-value} \\
		\midrule
		KMIS 2015 & 0.978 & (0.961-0.996) & 0.018 \\
		KMIS 2020 & 0.976 & (0.961-0.991) & 0.002 \\
		Combined & 0.972 & (0.960-0.984) & $<$0.001 \\
		\bottomrule
	\end{tabular}
\end{table}

\textbf{Interpretation of Table~\ref{tab:dml_ch4}:}

DML using XGBoost shows a consistent protective effect. For the combined dataset, OR = 0.972 (95\% CI: 0.960-0.984, p < 0.001), representing a 2.8\% reduction in odds.

This is smaller than the GLMM, PSM, and IPTW estimates (which showed 44-61\% reduction). The difference reflects:

\begin{enumerate}
	\item \textbf{DML estimates marginal effects}: Like DR, DML estimates the population-average treatment effect, not the cluster-specific effect.
	
	\item \textbf{Community effects are averaged in}: The DML estimate of 2.8\% mixes personal protection (where coverage is low) and community protection (where coverage is high). With cross-sectional data, we cannot separate these mechanisms.
	
	\item \textbf{Consistent with Hawley et al. (2003)}: Their cluster-randomized trial found that children in control compounds within 300 meters of ITN compounds received about the same protection as ITN users. The DML estimate likely captures the average of personal and community protection across Kenya.
\end{enumerate}

\subsection{E-value}

For the primary XGBoost estimate (OR = 0.972, 95\% CI: 0.960-0.984), the E-value is approximately 3.05.

\textbf{Interpretation:} An unmeasured confounder would need to be associated with both ITN use and malaria by a risk ratio of at least 3.05 to fully explain away the observed effect. For context, the strongest measured confounders (wealth, rainfall) have RRs in the range of 1.5-3.0. Thus, unmeasured confounding would need to be stronger than any measured confounder to eliminate the effect.

For the confidence interval bound closest to the null (OR = 0.986), the E-value is 1.86.

\section{Heterogeneity Analysis}

Table~\ref{tab:heterogeneity_ch4} presents subgroup-specific effects.

\begin{table}[htbp]
	\centering
	\caption{Heterogeneous Treatment Effects by Subgroup (DML)}
	\label{tab:heterogeneity_ch4}
	\begin{tabular}{@{}llccc@{}}
		\toprule
		\textbf{Subgroup} & \textbf{Category} & \textbf{N} & \textbf{OR} & \textbf{p-value} \\
		\midrule
		Wealth & Poorer & 1,448 & 0.919 & 0.001 \\
		Residence & Urban & 2,240 & 0.974 & 0.003 \\
		Age & 36-59 months & 3,075 & 0.963 & $<$0.001 \\
		Rainfall & Medium & 2,227 & 0.950 & $<$0.001 \\
		Rainfall & High & 2,156 & 0.967 & 0.025 \\
		\bottomrule
	\end{tabular}
\end{table}

\textbf{Interpretation of Table~\ref{tab:heterogeneity_ch4}:}

\textit{Wealth gradient:} ITNs are most effective for the "Poorer" quintile (OR = 0.919, p = 0.001) but not for the poorest (OR $\approx$ 1.00, not shown). This suggests the poorest households face additional barriers: poor housing quality (holes in walls, thatched roofs), overcrowded sleeping conditions (multiple children per net), or older, less effective nets due to inability to replace them. The "Poorer" quintile likely has better housing but still high baseline risk, giving ITNs more opportunity to demonstrate effect.

\textit{Residence gradient:} Urban areas show significant protection (OR = 0.974, p = 0.003); rural areas are borderline (OR = 0.976, p = 0.077). This may reflect differences in mosquito ecology (Anopheles gambiae bites indoors at night in both settings, but rural areas have higher vector density), housing quality (urban homes are more protective), and ITN usage patterns.

\textit{Age gradient:} Older children (36-59 months) show the strongest effect (OR = 0.963, p < 0.001). Younger children may kick off nets during sleep, and caregivers may not consistently cover them. Older children use nets more independently and consistently.

\textit{Rainfall gradient:} ITNs are most effective in medium and high rainfall areas (OR = 0.950-0.967). Transmission intensity is highest in these areas, giving ITNs more opportunity to demonstrate effect. In low rainfall areas, baseline malaria risk is already low, making it harder to detect a protective effect.

Figure~\ref{fig:heterogeneity} presents a forest plot of these heterogeneous effects.

\begin{figure}[htbp]
	\centering
	\includegraphics[width=0.8\textwidth]{Figures/week6_heterogeneity_forest.png}
	\caption{Heterogeneous Treatment Effects by Subgroup}
	\label{fig:heterogeneity}
\end{figure}

\textbf{Interpretation of Figure~\ref{fig:heterogeneity}:}

This forest plot displays odds ratios with 95\% confidence intervals for each subgroup.

- \textbf{X-axis}: Odds ratio (log scale)
- \textbf{Y-axis}: Subgroup category
- \textbf{Vertical dashed line at OR = 1}: Line of no effect
- \textbf{Points and bars}: Point estimates and 95\% confidence intervals

Key observations:

\begin{enumerate}
	\item \textbf{Wealth subgroups}: The "Poorer" quintile shows the strongest protective effect (OR = 0.919) with confidence interval entirely below 1. Other quintiles show effects closer to 1 with wider intervals.
	
	\item \textbf{Residence subgroups}: Urban shows significant protection; rural shows borderline protection with a wider interval.
	
	\item \textbf{Age subgroups}: The 36-59 month group shows the strongest effect with the narrowest interval (largest sample size, N = 3,075).
	
	\item \textbf{Rainfall subgroups}: Medium and high rainfall both show significant protection. Low rainfall shows no effect (OR $\approx$ 1.00).
\end{enumerate}

\section{Robustness Checks}

Table~\ref{tab:robustness_ch4} presents results from five alternative learners.

\begin{table}[htbp]
	\centering
	\caption{Robustness to Alternative Learners (Combined Dataset)}
	\label{tab:robustness_ch4}
	\begin{tabular}{@{}lccc@{}}
		\toprule
		\textbf{Learner} & \textbf{OR} & \textbf{95\% CI} & \textbf{p-value} \\
		\midrule
		RF (500 trees, node=10) & 0.977 & (0.957-0.996) & 0.021 \\
		RF (1000 trees, node=10) & 0.975 & (0.958-0.993) & 0.006 \\
		RF (500 trees, node=5) & 0.968 & (0.949-0.987) & 0.001 \\
		GB (100 rounds) & 0.980 & (0.969-0.991) & $<$0.001 \\
		GB (200 rounds) & 0.980 & (0.968-0.991) & $<$0.001 \\
		\bottomrule
	\end{tabular}
\end{table}

\textbf{Interpretation of Table~\ref{tab:robustness_ch4}:}

All five learners produced odds ratios between 0.968 and 0.980, all statistically significant. The narrow range (0.012 difference between the smallest and largest OR) demonstrates that DML estimates are robust to different machine learning specifications.

Figure~\ref{fig:robustness_ch4} compares the learners.

\begin{figure}[htbp]
	\centering
	\includegraphics[width=0.8\textwidth]{Figures/week6_robustness_comparison.png}
	\caption{Robustness of DML Estimates to Different Learners}
	\label{fig:robustness_ch4}
\end{figure}

\textbf{Interpretation of Figure~\ref{fig:robustness_ch4}:}

This forest plot shows ORs for each learner in descending order. All point estimates are below 1, and all confidence intervals exclude 1 (or nearly exclude it for the RF with node=10). The tight clustering of point estimates confirms robustness.

\section{Summary of Key Findings}

The main findings from this chapter are:

\begin{enumerate}
	\item \textbf{Confounding is substantial and negative}: Unadjusted models show no significant association. After demographic adjustment, ITN users appear at higher risk. Only after full environmental adjustment does the protective effect emerge (OR = 0.777 in combined sample, p = 0.049).
	
	\item \textbf{Clustering is substantial}: ICC = 50-58\%, justifying random effects models. The caterpillar plot shows substantial variation in cluster-specific malaria risk.
	
	\item \textbf{Conditional effects are large}: GLMM (conditional) shows OR = 0.557-0.588 (44-61\% reduction). PSM (conditional) shows OR = 0.392-0.466 (53-61\% reduction). IPTW (conditional) shows OR = 0.482-0.511 (48-52\% reduction).
	
	\item \textbf{Marginal effects are smaller}: DR shows near-null effects (OR $\approx$ 0.99). DML shows OR = 0.972 (2.8\% reduction). The difference between conditional and marginal effects is expected when ICC is high.
	
	\item \textbf{E-value = 3.05}: Unmeasured confounding would need to be quite strong (RR > 3) to explain away the effect.
	
	\item \textbf{Heterogeneity exists}: Strongest effects in poorer quintile, urban areas, older children, and high-rainfall areas.
	
	\item \textbf{Robustness confirmed}: All alternative learners produced ORs between 0.968-0.980.
\end{enumerate}